% sec-02: the intake sequence.  Figure 1 is a mermaid-type sequence diagram of
% the eleven-step institutional route recorded in priority-steps.md.
\section{The Route Being Followed}

\begin{appfloat}[!tb]
\begin{appfig}[mermaid-type / intake sequence]
% Five lifelines 2.9 apart; nine messages 0.72 apart.  Labels are set on a
% white ground so a message never overprints a lifeline.
\foreach \i/\x/\lab in {1/0/{Applicant}, 2/2.95/{Surgery and GI oncology},
  3/5.9/{Clinical Trials Office}, 4/8.85/{OCTA and ACTRI}, 5/11.8/{IRB}}{
  \node[mmactor] (a\i) at (\x,0) {\lab};
  \draw[mmlife] (\x,-0.42) -- (\x,-7.1);}
\foreach \y/\f/\t/\msg in {%
  -0.85/0/2.95/{1. Nonconfidential overview and feasibility request},
  -1.57/2.95/0/{2. Clinical supportability and site PI interest},
  -2.29/0/5.9/{3. Eligible-patient volume and OR capacity},
  -3.01/5.9/0/{4. Preliminary accrual estimate},
  -3.73/0/8.85/{5. CDA request and contracting-route decision},
  -4.45/8.85/0/{6. Sponsor-investigator determination},
  -5.17/0/5.9/{7. IIT concept form and budget questionnaire},
  -5.89/5.9/8.85/{8. Coverage analysis and budget},
  -6.61/8.85/11.8/{9. IRB submission, once and only once}}{
  \draw[mmmsg] (\f,\y) -- (\t,\y) node[midway,above,font=\tiny\sffamily,fill=protowhite,inner sep=1.2pt] {\msg};}
\node[font=\tiny\itshape,text=protogray,anchor=north west,text width=108mm]
  at (-0.6,-7.3) {Message 6 must precede message 9. Determining who can serve as
  sponsor-investigator after an IRB submission has been made forces the
  submission to be repeated.};
\end{appfig}
\vspace{-0.7cm}
\figcaption{The institutional route, message by message. The ordering
constraint that matters is the sponsor-investigator determination preceding the
IRB submission, not the total number of steps.}
\end{appfloat}
